% =====================================================================
%  pds-apuntes.tex — Estilo común de apuntes y hojas de ejercicios.
%
%  Uso (desde teoría/Tx/apuntes/ o teoría/Tx/ejercicios/):
%     \documentclass[11pt,a4paper]{scrartcl}   % o scrreprt
%     \input{../../comun/pds-apuntes}
%     \tema{1}{Introducción: señales y sistemas}
%     \begin{document} \cabecera{Apuntes} ... \end{document}
%  Compilar con:  latexmk -xelatex  (o «make» desde teoría/)
% =====================================================================

\input{../../comun/pds-comun.tex}

\usepackage[a4paper, top=2.6cm, bottom=2.6cm, left=2.5cm, right=2.5cm,
            headheight=1.1cm]{geometry}
\usepackage{fontspec}
\usepackage{unicode-math}
\setmainfont{LibertinusSerif}[Extension=.otf, UprightFont=*-Regular, BoldFont=*-Semibold,
  ItalicFont=*-Italic, BoldItalicFont=*-SemiboldItalic]
\setsansfont{FiraSans}[Extension=.otf, UprightFont=*-Regular, BoldFont=*-SemiBold,
  ItalicFont=*-Italic, BoldItalicFont=*-SemiBoldItalic, Scale=MatchLowercase]
\setmonofont{FiraMono}[Extension=.otf, UprightFont=*-Regular, BoldFont=*-Medium, AutoFakeSlant=0.15, Scale=0.85]
\setmathfont{LibertinusMath-Regular.otf}

\usepackage{microtype}
\usepackage{enumitem}
\usepackage{multicol}
\setlist{itemsep=2pt, topsep=4pt}
\usepackage{caption}
\captionsetup{font=small, labelfont={sf,bf,color=pdsazul}, labelsep=period}
\usepackage{subcaption}
\usepackage[most]{tcolorbox}
\usepackage{environ}
\usepackage{scrlayer-scrpage}
\usepackage[hidelinks, unicode]{hyperref}
\usepackage[spanish, noabbrev]{cleveref}

% Títulos en Fira Sans con el azul de la asignatura
\setkomafont{disposition}{\sffamily\bfseries\color{pdsazul}}
\setkomafont{descriptionlabel}{\sffamily\bfseries}
\setlength{\parindent}{0pt}
\setlength{\parskip}{0.55\baselineskip plus 2pt}
\linespread{1.05}

% ---------------------------------------------------------------------
%  Datos del tema y cabeceras
% ---------------------------------------------------------------------
\newcommand{\numtema}{}
\newcommand{\nombretema}{}
\newcommand{\tema}[2]{\renewcommand{\numtema}{#1}\renewcommand{\nombretema}{#2}}

\pagestyle{scrheadings}
\clearpairofpagestyles
\ihead{\sffamily\small\color{pdsgris}\asignatura}
\ohead{\sffamily\small\color{pdsgris}Tema~\numtema. \nombretema}
\ifoot{\sffamily\small\color{pdsgris}\grado}
\ofoot{\sffamily\small\pagemark}
\setkomafont{pagehead}{\normalfont}
\setkomafont{pagefoot}{\normalfont}

% \cabecera{Apuntes|Ejercicios|...} : bloque de título de la primera página
\newcommand{\cabecera}[1]{%
  \thispagestyle{plain.scrheadings}%
  \noindent
  \begin{minipage}[c]{2.2cm}
    \includegraphics[width=1.9cm]{../../comun/ucm.png}
  \end{minipage}%
  \begin{minipage}[c]{\dimexpr\linewidth-2.2cm}
    \sffamily
    {\small\color{pdsgris}\grado\par \asignatura\ · \curso\par}
    \vspace{4pt}
    {\Large\bfseries\color{pdsazul}#1\par}
    {\large\bfseries Tema~\numtema. \nombretema\par}
    \vspace{2pt}
    {\small\color{pdsgris}\autores\par}
  \end{minipage}\par
  \vspace{6pt}{\color{pdsazul}\rule{\linewidth}{0.8pt}}\par\vspace{4pt}}

% ---------------------------------------------------------------------
%  Cajas: definición, propiedad, ejemplo, nota, «en código»
% ---------------------------------------------------------------------
\tcbset{pdsbase/.style={enhanced, breakable, boxrule=0pt, arc=1.5pt, left=8pt, right=8pt,
    top=5pt, bottom=5pt, fonttitle=\sffamily\bfseries\small, coltitle=white,
    before skip=8pt, after skip=8pt}}
\newtcbtheorem[number within=section]{definicion}{Definición}%
  {pdsbase, colback=pdsfondo, colframe=pdsazul, borderline west={2.5pt}{0pt}{pdsazul}}{def}
\newtcbtheorem[use counter from=definicion]{propiedad}{Propiedad}%
  {pdsbase, colback=pdsfondo, colframe=pdsazul!75, borderline west={2.5pt}{0pt}{pdsazul!75}}{prop}
\newtcbtheorem[use counter from=definicion]{ejemplo}{Ejemplo}%
  {pdsbase, colback=pdsverde!5, colframe=pdsverde, borderline west={2.5pt}{0pt}{pdsverde}}{ej}
\newtcolorbox{nota}[1][Nota]{pdsbase, colback=pdsnaranja!7, colframe=pdsnaranja,
  title={#1}, borderline west={2.5pt}{0pt}{pdsnaranja}}
\newtcolorbox{importante}[1][Importante]{pdsbase, colback=pdsrojo!5, colframe=pdsrojo,
  title={#1}, borderline west={2.5pt}{0pt}{pdsrojo}}
\newtcolorbox{encodigo}[1][En código]{pdsbase, colback=white, colframe=pdsgris,
  title={#1}, borderline west={2.5pt}{0pt}{pdsgris}}

\lstset{basicstyle=\ttfamily\small}
\renewcommand{\codigoetiqueta}[1]{{\sffamily\small\bfseries\color{pdsgris}#1}}

% ---------------------------------------------------------------------
%  Hojas de ejercicios: \begin{ejercicio} ... \end{ejercicio} y
%  \begin{solucion} ... \end{solucion} (solo se imprime si se define
%  \SOLUCIONES antes de cargar el documento; ver *-Soluciones.tex)
% ---------------------------------------------------------------------
\newif\ifsoluciones
\ifdefined\SOLUCIONES \solucionestrue \fi
\newcounter{ejercicio}
\newenvironment{ejercicio}[1][]{%
  \refstepcounter{ejercicio}\par\vspace{6pt}%
  \noindent{\sffamily\bfseries\color{pdsazul}Ejercicio~\theejercicio}%
  \if\relax\detokenize{#1}\relax\else\ {\sffamily\color{pdsgris}(#1)}\fi\quad}%
  {\par}
\NewEnviron{solucion}{%
  \ifsoluciones
    \begin{tcolorbox}[pdsbase, colback=pdsverde!4, colframe=pdsverde,
      borderline west={2.5pt}{0pt}{pdsverde}, title={Solución}]
      \BODY
    \end{tcolorbox}
  \fi}
% \begin{incisos}[columnas] ... \end{incisos}: apartados a), b)... de un
% enunciado repartidos en columnas (3 por defecto)
% (con una sola columna no se usa multicols, que necesita al menos dos)
\newenvironment{incisos}[1][3]{\par\vspace{-2pt}%
  \ifnum#1>1 \def\finincisos{\end{multicols}}\begin{multicols}{#1}\raggedcolumns\raggedright
  \else\def\finincisos{}\fi
  \begin{enumerate}[label=\alph*), itemsep=3pt, topsep=0pt]}%
  {\end{enumerate}\finincisos\vspace{-4pt}}
\newcommand{\herramienta}[1]{\emph{Herramienta interactiva:} \texttt{herramientas/#1}}
